\chapter{Revisão bibliográfica e teoria da mudança}
\label{sec:literatura}

\section{Revisão bibliográfica}

\textbf{Agrotóxicos e desfechos perinatais.} A evidência causal vem de
desenhos que exploram onde e quando o produto é aplicado.
\citeonline{calzada2023bananas} estudam a fumigação aérea de bananais no
Equador com o endereço da mãe e o perímetro das plantações. Entre as mães que moram a até 150 metros, a
exposição intensa reduz o peso ao nascer em 80 gramas (erro-padrão de 31); o
efeito está no primeiro e no segundo trimestre de gestação e, além de 150
metros, os gradientes não se distinguem.
\citeonline{reynier2025glyphosate} encontram perda de 23 a 32 gramas e
gestação mais curta com a difusão do glifosato nos Estados Unidos.
\citeonline{larsen2017agricultural} só encontram efeito, de 5\% a 9\% nos
desfechos adversos, no topo da distribuição de exposição.
\citeonline{camacho2017health} associam a aspersão aérea de herbicida na
Colômbia a mais perdas gestacionais. O dano é localizado, concentra-se no
início da gestação e alcança a sobrevivência do feto.

\textbf{Evidência brasileira.} \citeonline{dias2023down} encontram aumento da
mortalidade infantil a jusante das áreas de maior uso de glifosato na soja.
\citeonline{januzzi2025} constrói um índice municipal de exposição a partir de
sete culturas: com efeitos fixos de município e de ano, os coeficientes para
baixo peso e prematuridade são positivos e imprecisos, e a associação com
baixo peso só aparece, a 10\%, na especificação sem efeitos de ano.
\citeonline{rodrigues2022tese} e \citeonline{rodrigues2022intoxicacoes}
relacionam a despesa com agrotóxicos por hectare às intoxicações notificadas,
que dependem da infraestrutura de vigilância. Essas medidas captam quantidade
ou gasto, e não método de aplicação, e não incluem a banana.

\textbf{A Chapada do Apodi.} \citeonline{rigotto2013trends} comparam Limoeiro
do Norte, Quixeré e Russas com municípios de agricultura familiar entre 2000 e
2010: a taxa de internação por neoplasia é 1,76 vez maior e a de óbito por
neoplasia, 1,38 vez; a razão de taxas não difere em óbito fetal nem em
malformação, embora a taxa de óbito fetal cresça nos municípios de
fruticultura. \citeonline{aguiar2017} descreve casos de malformação congênita
e puberdade precoce em uma comunidade rural. Desses estudos
descritivos o projeto toma o problema e o segundo grupo de comparação.

\textbf{Economia da regulação de agrotóxicos.} \citeonline{weitzman1974}
organiza a escolha entre regular por preço ou por quantidade. Do lado do
preço, \citeonline{bocker2016taxes} comparam os impostos sobre agrotóxicos de
quatro países europeus e concluem que o efeito é limitado, salvo
com alíquota alta, e \citeonline{finger2017taxation} defendem impostos
diferenciados pelo risco de cada produto. No Brasil o instrumento de preço
opera como subsídio: \citeonline{cunha2020incentivos} calculam em R\$~9,86
bilhões a desoneração fiscal dos agrotóxicos em 2017.
E \citeonline{seixas2022gm} mostram que uma mudança de tecnologia pode elevar o
impacto ambiental por substituição entre produtos: convém medir, e não supor,
o que substituiu o avião.

\section{Teoria da mudança}

A lógica da lei está escrita na justificativa do projeto que lhe deu origem
\cite{ceara2015pl18} e foi retomada no parecer técnico entregue à Assembleia
\cite{tramas2016parecer} e no julgamento do Supremo \cite{stf2023adi6137}. O
problema é a deriva, que alcança comunidades e mananciais; o fundamento, os
princípios da precaução e da prevenção. A cadeia é: \emph{se} a lei proíbe a aplicação por aeronave e a proibição é
cumprida, \emph{então} cessa a deriva, cai a exposição de gestantes e
melhoram os desfechos perinatais.

Três pontos importam para a avaliação. A premissa de que a prática estava em
curso não foi medida antes da lei. O instrumento proíbe o método e não a molécula: os mesmos produtos seguem
permitidos por via terrestre, e a troca pode reduzir a exposição ambiental e
elevar a ocupacional. E o funil é estreito: a prática precisava existir, a
proibição ser cumprida, a mãe morar ao alcance da deriva e a gestação
atravessar a estação de aplicação.

A Figura~\ref{fig:dag} representa a cadeia como grafo causal. Uma avaliação de
impacto estima o efeito total da lei, por todos os caminhos: o caminho-alvo; a
substituição por aplicação terrestre, que também expõe vizinhos e
trabalhadores; a renda; a composição das mães; e a seleção para nascer vivo,
que atenua o peso médio se o óbito fetal cair. A análise dos demais caminhos é
evidência de implementação, e não decompõe o efeito total.

\begin{figure}[htb]
\centering
\caption{Grafo causal resumido}
\label{fig:dag}
\begin{tikzpicture}[>=Stealth, node distance=4mm and 13mm,
  every node/.style={font=\footnotesize},
  obs/.style={draw, rounded corners=2pt, align=center, inner sep=3pt},
  lat/.style={draw, dashed, rounded corners=2pt, align=center, inner sep=3pt}]
  \node[obs] (L) {Lei};
  \node[obs, right=of L] (A) {Aplicação\\aérea};
  \node[lat, right=of A] (X) {Exposição\\ambiental};
  \node[obs, right=of X] (Y) {Desfecho\\ao nascer};
  \node[obs, below=of A] (T) {Aplicação\\terrestre};
  \node[lat, below=of X] (O) {Exposição\\ocupacional};
  \node[obs, above=of A] (P) {Produção\\e emprego};
  \node[obs, above=of P] (C) {Composição\\das mães};
  \node[obs, below=of Y] (F) {Perda\\fetal};
  \node[obs, right=of Y] (G) {Choque\\estadual};
  \draw[->, very thick] (L) -- (A);
  \draw[->, very thick] (A) -- (X);
  \draw[->, very thick] (X) -- (Y);
  \draw[->] (L) |- (T);
  \draw[->] (T) -- (X);
  \draw[->] (T) -- (O);
  \draw[->] (O) -- (Y);
  \draw[->] ([xshift=1.5mm]L.north) |- (P);
  \draw[->] ([xshift=-1.5mm]L.north) |- (C);
  \draw[->] (P) -| ([xshift=-2mm]Y.north);
  \draw[->] (C) -| ([xshift=2mm]Y.north);
  \draw[->] (X) -- (F);
  \draw[->] (F) -- (Y);
  \draw[->, dashed] (G) -- (Y);
\end{tikzpicture}
\fonte{Elaboração própria. Em traço grosso, o caminho-alvo. Caixas tracejadas: variáveis não observadas. Seta tracejada: ameaça à identificação. A perda fetal afeta o desfecho observado por seleção de quem nasce vivo.}
\end{figure}
